\documentclass[oneside]{report}
\usepackage{phd}

\definecolor{spot}{rgb}{0,0.2,0.6}
\definecolor{spot}{rgb}{0.95,.48,.38}


\definecolor{theblue} {rgb}{0.02,0.04,0.48}
\definecolor{thegrey} {gray}{0.5}
\definecolor{thegray} {gray}{0.5}
\def\ZZZ{\phantom{ZZZ}}
\makeatletter
\def\HHUGE{\@setfontsize\HHUGE{58}{89}\selectfont}
\makeatother
\usepackage{tikz}
\usetikzlibrary{%
  calc,%
  backgrounds,%
  fadings,%
  positioning,%
  tikzmark
   }

\usepackage[bottom=2.5cm]{geometry}
\fancyhf{}
\renewcommand{\headrulewidth}{0pt}
\pagestyle{fancy}
\thispagestyle{fancy}
\lhead[E]{ %
 \begin{tikzpicture}[remember picture, overlay]
 \node [xshift=0cm,yshift=-1.5cm] at (current page.north west)
[text width=3cm, text=white ,text height=2cm, shade, left color=spot!50!white, right color=white, above right]
{
\hbox to 1cm{\hfill \textbf{\thepage}}
\vspace*{1em}
};
 \end{tikzpicture}
 \color{white}}

\begin{document}

\makebox[\textwidth]{\hfill
\begin{minipage}{.8\textwidth}
\columnseprule2pt
\def\columnseprulecolor{\color{thegray}}
\columnsep22pt
\begin{multicols}{2}

\color{theblue}
\flushright
\Large
\obeylines
Over the last year we
have continued to
develop and improve the
range of funding schemes
we offer to meet the
needs of the arts and
humanities communities,
for example, by offering
opportunities for early
career researchers.
\columnbreak
\color{thegray}

\small
\flushleft

\obeylines
We have engaged both
individuals and groups to
build a vision for our strategic
initiatives and our museums
and galleries strategy, have
opened up opportunities
for the arts and humanities
in cross-Council funding
initiatives and undertaken
to represent the needs of our
communities in arenas such
as the Research Councils’
project on the Efficiency and
Effectiveness of Peer Review
Journals
initiatives and undertaken
to represent the needs of our
communities in arenas such
as the Research Councils’
project on the Efficiency and
Effectiveness of Peer Review
Journals
\end{multicols}
\end{minipage}
}
\bigskip
\bgroup


\fboxsep20pt\color{thegray!60}
\hspace*{-1.85cm}\fbox{\HHUGE Research}

\bigskip
\egroup

\begin{multicols}{2}
\color{black!80}

\parindent1em
\lettrine{L}{ipsum} \lipsum[1-3]
\end{multicols}

\newpage

\hspace*{-2.2cm}\begin{minipage}{\textwidth+2cm}
\columnseprule2pt
\def\columnseprulecolor{\color{spot!50!white}}
\columnsep22pt
\begin{multicols}{2}
\color{spot}
\flushright
\LARGE
\obeylines
One of the challenges in
the area of knowledge
transfer is to identify
models of opportunity
for engagement.
\columnbreak
\color{thegray}


\large
\flushleft

\advance\leftskip-0.2em %

Opportunities that build trust between\\
researchers and other partners; enable\\
an exchange of knowledge and ideas that\\
benefit both; and encourage novel forms\\
of engagement that permit academic and\\
non-academic partners to maximize their\\
contribution to the nation’s cultural, social\\
and economic life. \\
\endgraf
\rule{0cm}{1cm}
\end{multicols}
\end{minipage}
\vskip25pt

\tikzmark{p0}

 {\fboxsep4pt\fboxrule0pt\color{spot!50!white}
  \hskip\dimexpr-\leftmargin-\hoffset\relax\fbox{\HHUGE\color{spot}\tikzmark{p3}Production \tikzmark{p1}}
  \vskip3pt
  \makebox[\textwidth]{\HHUGE \color{spot}\hfill\tikzmark{p4} \hbox{Technologies} \tikzmark{p2}}
  }

  \tikzmark{p6}
 \begin{tikzpicture}[overlay, remember picture, color=spot!50, line width=2pt]
    \draw({pic cs:p0} -| current page.west) +(0,5pt) -| ([yshift=-0pt,xshift=5pt]{pic cs:p1}) -| ([xshift=5pt]{pic cs:p2})  |- ({pic cs:p4} |- {pic cs:p6}) -| ([yshift=-5pt, xshift=-5pt]{pic cs:p4})[] |- ([yshift=-20pt]{pic cs:p3}) coordinate (p7) -- (p7 -| current page.west);
  \end{tikzpicture}


\bigskip
\parindent1em


\begin{multicols}{2}

\lettrine{P}{ost}  \lipsum
\end{multicols}





  \parindent1em

 
\makeatletter
The value of $x$\pgfsys@markposition{here} is important.
Lots of text.
\hbox{\pgfsys@markposition{myorigin}%
\begin{pgfpicture}
% Switch of size protocol
\pgfpathmoveto{\pgfpointorigin}
\pgfusepath{use as bounding box}
\pgfsys@getposition{here}{\hereposition}
\pgfsys@getposition{myorigin}{\thispictureposition}
\pgftransformshift{\pgfpointscale{-1}{\thispictureposition}}
\pgftransformshift{\hereposition}
\pgfpathcircle{\pgfpointorigin}{0.5cm}
\pgfusepath{draw}
\end{pgfpicture}}

The header and footer are constructed using three parboxes. One for left, one for center and one for right

\def\@fancyhead#1#2#3#4#5{\raggedright%
   #1\hbox to\headwidth{\fancy@reset
  \@fancyvbox\headheight{\headrule\hbox
    {
     \rlap{\fbox{\parbox[t]{\headwidth}{\raggedright#2}}}\hfill
      \fbox{\parbox[t]{\headwidth}{\centering#3}}\hfill
      \llap{\fbox{\parbox[t]{\headwidth}{\raggedleft#4}}}}\headrule}}#5}

\def\headrulewidth{2pt}
\def\headwidth{210mm}
\@fancyhead{\hspace*{-3.1cm}}{left longish text to see\\ what is happening here}{\tikzmark{t11}center mark\\other\\maybe\\all\\ \begin{tabular}{ll}1&2\\\end{tabular}}{right text }{\thepage \tikzmark{t12}}

 \begin{tikzpicture}[overlay, remember picture, color=spot!50!white, line width=\fboxrule]
    \draw (pic cs:t11) -- (pic cs:t12) -- (pic cs:p1)[->,thick] ;
  \end{tikzpicture}

\makeatother














\end{document}